\newpage
\datos{color}{}{}
\section*{Propuesta técnica de búsqueda evolutiva}
\addcontentsline{toc}{section}{Propuesta técnica de búsqueda evolutiva}

\subsection*{Algoritmo elegido}
Se ha seleccionado la técnica de Evolución Diferencial (\textit{DE}). 
Esta elección se justifica por la naturaleza del problema definido en los Temas 1 y 2, 
que opera en un espacio de búsqueda continuo con variables 
reales. Se trabajará sobre el planteamiento derivado de la revisión del Tema 2.
Cabe destacar que al principio, dadas las restricciones del problema, especialmente 
la relación espacial de los puntos de encuentro, se planteó el algoritmo 
de Nubes de Partículas (\textit{PSO}). Sin embargo, tras la tutorización del equipo
docente y a la vista de los resultados obtenidos en el Tema 2, se determinó 
que la \textit{DE} tenderá a mantener una mayor diversidad poblacional 
y ser menos propensa a la convergencia prematura. 

%Frente a los Algoritmos Genéticos (\textit{GA}), la Evolución Diferencial
%ofrece una mayor naturalidad para trabajar con codificación real sin necesidad de diseñar operadores 
%de cruce y mutación complejos para variables heterogéneas. \\ Por su parte, en comparación con el algoritmo 
%de Nubes de Partículas (\textit{PSO}), la \textit{DE} tiende a mantener una mayor diversidad poblacional 
%y es menos propensa a la convergencia prematura. 

\subsection*{Diseño del algoritmo}
\subsubsection*{Representación de las soluciones}
Cada solución candidata se representa como el vector continuo de la \autoref{vector}.
%\[
%\mathbf{x} =
%(\mathbf{p}_A,\mathbf{p}_B,
%\mathbf{q}_A,\mathbf{q}_B,
%s_A,s_B,s_{rail})
%\in \mathbb{R}^{18},
%\]
%manteniendo todas las especificaciones declaradas en la revisión del segundo tema.

\subsubsection*{Planteamiento del problema}
Se continúa con el planteamiento mono-objetivo penalizado que se ha desarrollado en la 
revisión del Tema 2, adoptando una
arquitectura de dos niveles para reducir la dimensionalidad y filtrar soluciones
cinemáticamente inválidas antes de la optimización. Al igual que con la búsqueda local,
el algoritmo evolutivo se aplica al Nivel 2, donde el vector de decisión se centra en las 
variables continuas que regulan la ejecución temporal y la posición final del raíl: 
\begin{equation}
    \mathbf{x} = (s_A, s_B, s_{rail}, p_{rail}) \in \mathbb{R}^4
\end{equation}

\subsubsection*{Operadores evolutivos}
Se utilizará la variante estándar \textbf{DE/rand/1/bin}:
\begin{itemize}
    \item \textbf{Inicialización:} población generada aleatoriamente mediante una 
distribución uniforme en el espacio de búsqueda $(0, 1]$ para garantizar diversidad inicial. 
El rango de las variables de velocidad es $(0, 1]$, por lo que se proyecta el rango real de la
posición del raíl a ese mismo intervalo para facilitar los mecanismos del algoritmo.
    \item \textbf{Perturbación:} se genera un vector mutante mediante la 
diferencia escalada de dos individuos aleatorios sumada a un tercero: $v_i = x_{r1} + F \cdot (x_{r2} - x_{r3})$.
    \item \textbf{Cruce:} cruce binomial (uniforme) que combina el vector 
actual con el mutante basado en una probabilidad $CR$.
    \item \textbf{Selección:} selección elitista 1 a 1, donde el hijo 
sustituye al padre solo si mejora su valor de aptitud o factibilidad.
    \item \textbf{Reparación:} tras el cruce, se aplica una función de reparación 
para proyectar la posición del raíl obtenida a sus límites físicos permitidos.
\end{itemize}

\subsubsection*{Método de evaluación}
Se va a emplear el mismo tratamiento de restricciones que en el Tema 2. 

%La función de mérito evalúa el coste temporal $T_{max}$ e integra penalizaciones 
%continuas para la desincronización $|t_A - t_B|$ \cite{12}. Para las restricciones 
%críticas (cinemática inversa y colisiones), se aplican funciones de penalización que 
%aumentan el coste proporcionalmente a la violación de la restricción, guiando al algoritmo
%hacia regiones factibles \cite{12, 15}.

\subsection*{Configuración de parámetros}
Siguiendo las indicaciones del equipo docente y las recomendaciones técnicas para optimización en espacios continuos, se propone la siguiente configuración:

\begin{itemize}
    \item \textbf{Tamaño de la población ($NP$) de 60:} se justifica bajo la regla 
heurística de emplear entre $10d$ y $20d$, siendo $d = 4$ la dimensionalidad del \textit{DE}.

    \item \textbf{Factor de escala ($F$) de 0.75:} situado en el rango recomendado de 
$[0.5, 0.8]$, permite que las perturbaciones diferenciales basadas en la 
diversidad de la población sean lo suficientemente grandes para saltar óptimos locales.

    \item \textbf{Probabilidad de cruce ($CR$) de 0.90:} se selecciona un valor alto 
debido al fuerte acoplamiento entre las velocidades de los robots y la posición del raíl. 
Esto asegura que el vector hijo herede la mayor parte de la información del vector mutante
perturbado, evitando el estancamiento en regiones subóptimas.

    \item \textbf{Máximo de generaciones (150):} este valor se ha fijado para igualar 
el coste computacional del Tema 2. Dado que en el Temple Simulado se realizaron 300 
iteraciones para cada uno de los 30 puntos espaciales (9,000 evaluaciones totales), 
establecer 150 generaciones con una población de 60 individuos genera exactamente el mismo 
número de soluciones ($60 \cdot 150 = 9,000$).
\end{itemize}

Además se proponen tres grupos más de parámetros para realizar un análisis de 
sensibilidad de \textit{Kruskal-Wallis}: una configuración exploratoria
($F=0.4$, $CR=0.5$), una intermedia ($F=0.6$, $CR=0.7$) y una con perturbaciones agresivas ($F=0.9,CR=0.9$).

\subsection*{Diseño experimental}
El plan de experimentación consistirá en 30 ejecuciones independientes por 
cada grupo de parámetros a evaluar mediante el test estadístico de \textit{Kruskal-Wallis}. 
Finalmente, se 
seleccionará la configuración más adecuada y sus resultados se contrastarán
con los obtenidos en el tema anterior, utilizando para ello los mismos diagramas, esquemas y 
tablas. 